\begin{table*}[t!]
    \centering
    \scalebox{0.85}{
    \begin{tabular}{cl}
        \toprule
        Mode & Prompt \\
        \midrule
        Default & \texttt{\thead{You are an expert programmer. What does the following code snippet in\\Python 3.7 print?\\\textasciigrave\textasciigrave\textasciigrave python\\def function(lst):\\\ \ \ \ return sum([lst[i] for i in range(1, len(lst), 2) if lst[i] \% 2 == 0])\\\\print([function([4, 88])])\\print([function([4, 5, 6, 7, 2, 122])])\\print([function([4, 0, 6, 7])])\\print([function([4, 4, 6, 8])])\\print([list(range(3))])\\print([[4, 5, 6].pop(2)])\\print(["qrs"[:2]])\\print(["qrstu"[4]])\\print([list(enumerate("qrstuv"))])\\\textasciigrave\textasciigrave\textasciigrave\\\textbf{\{Let's think step by step. Write out intermediate results and reasoning}\\\textbf{processes as needed. \}}End the response by saying "The final output is:"\\and a unified summary \textasciigrave\textasciigrave\textasciigrave python\textasciigrave\textasciigrave\textasciigrave\ code block with *ALL* the output, in\\which each line represents the output of each print statement.}} \\
        \midrule
        CF & \texttt{\thead{You are an expert programmer who can readily adapt to new programming\\languages. There is a new programming language, ThonPy, which is identical to\\Python 3.7 except all variables of the \textasciigrave list\textasciigrave , \textasciigrave tuple\textasciigrave , and \textasciigrave str\textasciigrave\ types use\\1-based indexing, like in the MATLAB and R languages, where sequence indices\\start from 1. That is, index \textasciigrave n\textasciigrave\ represents the \textasciigrave n\textasciigrave-th element in a sequence,\\NOT the \textasciigrave n+1\textasciigrave-th as in 0-based indexing. This change only affects when the\\index is non-negative. When the index is negative, the behavior is the same\\as Python 3.7. This also affects methods of these classes such as \textasciigrave index\textasciigrave\ and\\\textasciigrave pop\textasciigrave. The built-in functions \textasciigrave enumerate\textasciigrave\ and \textasciigrave range\textasciigrave\ also use 1-based\\indexing: by default, the index of \textasciigrave enumerate\textasciigrave\ starts from 1, and so does the\\lower bound of \textasciigrave range\textasciigrave\ when not supplied (the higher bound is unchanged).\\\\For example,\\\textasciigrave\textasciigrave\textasciigrave thonpy\\assert (7, 8, 9)[1] == 7\\assert ["abc", "def", "ghi"][3] == "ghi"\\assert "abcde"[4] == "d"\\assert "abc"[:2] == "a"\\assert [7, 8, 9][1:] == [7, 8, 9][1:5] == [7, 8, 9][1::1] == [7, 8, 9][:4]\\== [9, 8, 7][::-1] == [9, 8, 7, 6][3::-1] == [7, 8, 9]\\assert list(enumerate([7, 8, 9])) == [(1, 7), (2, 8), (3, 9)]\\assert list(range(2)) == [1]\\assert list(range(2, 4)) == [2, 3]\\assert \{0: 7, 1: 8, 2: 9\}[1] == 8\\assert [7, 8, 9].index(8) == 2\\\textasciigrave\textasciigrave\textasciigrave\\\\What does the following code snippet in ThonPy print?\\\textasciigrave\textasciigrave\textasciigrave thonpy\\def function(lst):\\\ \ \ \ return sum([lst[i] for i in range(1, len(lst), 2) if lst[i] \% 2 == 0])\\\\print([function([4, 88])])\\print([function([4, 5, 6, 7, 2, 122])])\\print([function([4, 0, 6, 7])])\\print([function([4, 4, 6, 8])])\\print([list(range(3))])\\print([[4, 5, 6].pop(2)])\\print(["qrs"[:2]])\\print(["qrstu"[4]])\\print([list(enumerate("qrstuv"))])\\\textasciigrave\textasciigrave\textasciigrave\\\textbf{\{Let's think step by step. Write out intermediate results and reasoning}\\\textbf{processes as needed. \}}End the response by saying "The final output is:"\\and a unified summary \textasciigrave\textasciigrave\textasciigrave thonpy\textasciigrave\textasciigrave\textasciigrave\ code block with *ALL* the output, in\\which each line represents the output of each print statement.}} \\
        \bottomrule
    \end{tabular}
    }
    \caption{\label{tab:prompts-programming}
    Prompts for the program execution task. \texttt{\textbf{\{Let's think step by step. Write out intermediate results and reasoning processes as needed. \}}} is added only if 0-shot CoT is used. All the print statements wrap the expression in a singleton list for the ease of parsing, so that (a) each output always takes a single line even with line breaks in the middle, and (b) we can distinguish between a string representation of e.g. an integer and the integer type.
    }
\end{table*}



\begin{table*}[t!]
    \centering
    \begin{tabular}{cl}
        \toprule
        Mode & Prompt \\
        \midrule
        Default & \texttt{\thead{You are an expert programmer. Complete the following function in Python 3.7. Please\\only output the code for the completed function.\\\\\\def add(lst):\\\ \ \ \ """Given a non-empty list of integers lst. add the even elements that are at odd\\indices..\\\\\\\ \ \ \ Examples:\\\ \ \ \ \ \ \ \ add([4, 2, 6, 7]) ==> 2 \\\ \ \ \ """\\}} \\
        \midrule
        CF & \texttt{\thead{You are an expert programmer who can readily adapt to new programming\\languages. There is a new programming language, ThonPy, which is identical to\\Python 3.7 except all variables of the \textasciigrave list\textasciigrave , \textasciigrave tuple\textasciigrave , and \textasciigrave str\textasciigrave\ types use\\1-based indexing, like in the MATLAB and R languages, where sequence indices\\start from 1. That is, index \textasciigrave n\textasciigrave\ represents the \textasciigrave n\textasciigrave-th element in a sequence,\\NOT the \textasciigrave n+1\textasciigrave-th as in 0-based indexing. This change only affects when the\\index is non-negative. When the index is negative, the behavior is the same\\as Python 3.7. This also affects methods of these classes such as \textasciigrave index\textasciigrave\ and\\\textasciigrave pop\textasciigrave. The built-in functions \textasciigrave enumerate\textasciigrave\ and \textasciigrave range\textasciigrave\ also use 1-based\\indexing: by default, the index of \textasciigrave enumerate\textasciigrave\ starts from 1, and so does the\\lower bound of \textasciigrave range\textasciigrave\ when not supplied (the higher bound is unchanged).\\\\For example,\\\textasciigrave\textasciigrave\textasciigrave thonpy\\assert (7, 8, 9)[1] == 7\\assert ["abc", "def", "ghi"][3] == "ghi"\\assert "abcde"[4] == "d"\\assert "abc"[:2] == "a"\\assert [7, 8, 9][1:] == [7, 8, 9][1:5] == [7, 8, 9][1::1] == [7, 8, 9][:4]\\== [9, 8, 7][::-1] == [9, 8, 7, 6][3::-1] == [7, 8, 9]\\assert list(enumerate([7, 8, 9])) == [(1, 7), (2, 8), (3, 9)]\\assert list(range(2)) == [1]\\assert list(range(2, 4)) == [2, 3]\\assert \{0: 7, 1: 8, 2: 9\}[1] == 8\\assert [7, 8, 9].index(8) == 2\\\textasciigrave\textasciigrave\textasciigrave\\\\Complete the following function in ThonPy. Please only output the code for the\\completed function.\\\\\\def add(lst):\\\ \ \ \ """Given a non-empty list of integers lst. add the even elements that are at odd\\indices..\\\\\\\ \ \ \ Examples:\\\ \ \ \ \ \ \ \ add([4, 2, 6, 7]) ==> 2 \\\ \ \ \ """\\}} \\
        \bottomrule
    \end{tabular}
    \caption{\label{tab:prompts-programming-gen}
    Prompts for the program generation task.
    }
\end{table*}
